\documentclass[11pt,oneside]{book}

\usepackage[T1]{fontenc}
\usepackage[utf8]{inputenc}
\usepackage{amsmath,amssymb,mathtools}
\usepackage{booktabs}
\usepackage[margin=1in]{geometry}
\usepackage{xcolor}
\usepackage{listings}
\usepackage[hidelinks]{hyperref}

\definecolor{codeblue}{RGB}{25,75,135}
\definecolor{codegreen}{RGB}{25,115,65}
\definecolor{codegray}{RGB}{95,95,95}
\definecolor{codebackground}{RGB}{247,247,247}

\lstdefinestyle{python}{
  language=Python,
  basicstyle=\ttfamily\small,
  keywordstyle=\color{codeblue}\bfseries,
  commentstyle=\color{codegreen},
  stringstyle=\color{codegray},
  backgroundcolor=\color{codebackground},
  frame=single,
  rulecolor=\color{black!20},
  numbers=left,
  numberstyle=\tiny\color{codegray},
  numbersep=8pt,
  showstringspaces=false,
  breaklines=true,
  breakatwhitespace=false,
  columns=fullflexible,
  keepspaces=true,
  tabsize=4
}
\lstset{style=python}

\newcommand{\Hubbard}{\mathcal{H}}
\newcommand{\basisdim}{\mathcal{D}}
\newcommand{\ket}[1]{\left\lvert #1 \right\rangle}
\newcommand{\bra}[1]{\left\langle #1 \right\rvert}
\newcommand{\braket}[2]{\left\langle #1 \middle\vert #2 \right\rangle}
\newcommand{\expect}[1]{\left\langle #1 \right\rangle}
\newcommand{\norm}[1]{\left\lVert #1 \right\rVert}
\newcommand{\code}[1]{\texttt{#1}}

\hypersetup{
  pdftitle={Building a Trustworthy Hubbard Exact-Diagonalization Solver},
  pdfauthor={hubbard-ed contributors},
  pdfsubject={Exact diagonalization tutorial for the one-dimensional Hubbard model}
}
